\section{Conclusion}

We set out to test whether output diversity can be controlled by adding a direction to the residual
stream, estimated by contrasting a model's more and less varied generations. The answer is no, and
the reason is more interesting than the answer. The estimator appears to work when checked the obvious way:
two halves yield the same direction to three decimals on pythia-410m. It
stops working the moment both halves come from the same prompt, where the cosine falls to $-0.037$
in that model and never exceeds $|0.12|$ across our seven models. The pooled statistic recovers prompt identity, which is reproducible for reasons unrelated to
dispersion. The pooled vector is consequently
a sum of near-orthogonal prompt-local components, none of which predicts the others, which is why
injecting it across thirty-five dose levels produces no positive dose-response in any model, why it
does not outperform a random vector of matched norm, and why the one prompt-specific variant we
tested performs no better than borrowing another prompt's direction.

Ablation tells a different and more useful story. Removing the response subspace collapses diversity
in pythia-1b from $6.13$ to $1.35$ at layer 4, and a control that keeps the activations but destroys
the dispersion pairing collapses it to $1.46$, so the causal factor is the subspace and not the
contrast. Diversity is thus one-sided: it can be destroyed by removing the span of a prompt's
responses, but it cannot be increased by adding a direction, because there is none to add.
We hope these four checks become routine; each is cheap and each rules out a class of result that
would otherwise look like success.
